\section{Introduction}\label{sec:intro}

Let $X$ be a smooth complex projective variety. The Hodge conjecture asserts
that the cycle class map
\[
  \cl^{p}_{X}\colon \CH^{p}(X)_{\QQ}\longrightarrow \Hdg^{p}(X)
  \coloneqq H^{2p}(X,\QQ)\cap H^{p,p}(X)
\]
is surjective for every $p$. Beyond degree two the best understood Hodge
classes are those of abelian varieties, and the first of them that divisors
do not produce are the Weil classes, which Weil exhibited on abelian
varieties with an action of an imaginary quadratic field whose two
eigenspaces on the tangent space have the same dimension \cite{Weil77}; the
construction works for every CM field $F$. Markman proved that the Weil
classes are algebraic on the abelian sixfolds of split Weil type for every
imaginary quadratic field, and, through a product argument of Schoen
\cite[Proposition~10]{Sch98}, on every abelian fourfold of Weil type
\cite{Mar25a,Mar25b}. His sheaves are built from secant sheaves on an abelian
variety of half the dimension, and they extend along the family because they
satisfy an equivariant form of semiregularity.

The companion paper \cite{paperI} studies the set of members of a family on
which a flat Hodge class is algebraic, its \emph{algebraic locus}, and
determines it for the Weil classes of every CM field and every dimension at
once. Fix a CM field $F$ of degree $2m$, an integer $n\ge2$ and a
discriminant $\delta$. The abelian varieties of $(F,n,\delta)$-Weil type are
the members of one family over a connected period domain $\cD_{F}$ of
dimension $mn^{2}$, their Weil classes form a flat local system of rank $2m$
over it, and the algebraic locus
\[
  \Sigma_{F}=\{\,s\in\cD_{F}:\text{the Weil classes of }A_{s}\text{ are
  algebraic}\,\}
\]
contains an explicit point of complex multiplication, is dense, is the
preimage of a countable union of Zariski closed subsets of the Shimura
variety of the family, and is either all of $\cD_{F}$ or meagre; so the
following are equivalent: $\Sigma_{F}$ is closed, it has an interior point,
the Weil classes are algebraic on every member, and the Hodge conjecture
holds for the very general member (\cite[Theorem~1.1]{paperI}, which we call the main
theorem throughout). The notation of that paper is kept here: the standard
model of \cite[Section~4]{paperI}, the integral basis $\omega_{1},\omega_{2}$ of
the Weil line of \cite[Theorem~4.9]{paperI}, the base points of
\cite[Theorem~15.27]{paperI} and \cite[Theorem~17.5]{paperI}, and the strata of
\cite[Theorem~16.4]{paperI}.

This paper determines what closing $\Sigma_{F}$ requires of a single object
at a single point, computes that requirement exactly, and examines the
objects that the known constructions supply, showing in each case where they
fail. The third paper \cite{paperIII} places the Weil classes inside the
whole conjecture.

\subsection{Closing the algebraic locus}

By \cite[Theorem~1.1(iv)]{paperI} the Hodge conjecture for the very general member of a
Weil family is the closedness of $\Sigma_{F}$, and by (iii) it is enough to
show that one stratum has an interior point. Two statements would do this,
and the paper determines what each of them asks.

The first bounds the Hilbert data. If at every point of the orbit of the
base point the Weil class is represented by a cycle whose Hilbert data comes
from one finite set, then finitely many strata cover a dense set and one of
them is everything (\cite[Theorem~16.9]{paperI}). Read on an arbitrary
representative this bound is equivalent to its conclusion
(\cite[Theorem~16.17]{paperI}); read on the transported cycle, the one
representative that can be named, it fails, because the transport multiplies
the degree of a component by $c^{2n}$ and divides it only by the order of a
stabiliser (\cite[Theorem~16.15]{paperI}). So the bound restates the problem rather
than reducing it.
The bound is needed only on a small set: by the theorem of Pila, Shankar
and Tsimerman on the Andr\'e--Oort conjecture \cite{PST},%
\footnote{In the form used here: an irreducible closed subvariety of a
Shimura variety of abelian type that contains a Zariski dense set of CM
points is a special subvariety. Hence a sequence of CM points that leaves
every proper special subvariety is Zariski dense, and so is every infinite
subsequence of it.}
it suffices along
one sequence of CM points that leaves every proper special subvariety, and
such a sequence exists inside the orbit (\cite[Theorem~17.14]{paperI}). Every
member at which this paper exhibits algebraic Weil classes has endomorphisms
beyond the field, so it lies in the locus of members with smaller Hodge
group. That locus contains the base point, is dense, is stable under
$G_{F}(\QQ)$ and is a countable union of closed algebraic strata, exactly as
$\Sigma_{F}$ is, and it is meagre (\cite[Proposition~17.12]{paperI},
\cite[Corollary~17.13]{paperI}). The properties listed in \cite[Theorem~1.1]{paperI} therefore
do not decide the dichotomy by themselves, and a proof that $\Sigma_{F}$ is
everything has to reach a member with full Hodge group.

The second is semiregularity. By the theorem of Buchweitz and Flenner, in the
form of Pridham for perfect complexes \cite{BF03,Pri24},%
\footnote{The semiregularity map of a perfect complex $E$ on $X$ sends
$\Ext^{2}(E,E)$ to $\bigoplus_{q}H^{q+2}(X,\Omega^{q}_{X})$ by the trace of
the product with powers of the Atiyah class of $E$; it carries the
obstruction to deforming $E$ along a first order deformation of $X$ to the
obstruction to keeping $\ch(E)$ of Hodge type. If it is injective, $E$
deforms, formally and then along the whole family, in every direction in
which $\ch(E)$ stays a Hodge class. Bloch introduced it for local complete
intersection subvarieties.}
a perfect complex at
one base point whose Chern character is a nonzero multiple of a Weil class
plus a polynomial in the polarisation, and whose semiregularity map is
injective, extends along the whole family and makes $\Sigma_{F}$ everything
(\Cref{thm:semiregclosure}, \cite[Theorem~17.7(v)]{paperI}). The paper computes
the injectivity condition exactly. The annihilator of the Weil class in
$H^{0,2}$ is carried by the polarisation isomorphically onto the tangent space
of the family (\cite[Lemma~9.6]{paperI}, \cite[Theorem~9.17]{paperI}), so no object
whose Chern character lies on the Weil line and on which $H^{2}(\cO_{A})$ acts
faithfully is semiregular (\cite[Theorem~9.7]{paperI}). When the Chern character is
exactly a Weil class the criterion is never met, for any CM field and any
$n\ge2$: a semiregular object would extend to the very general Weil torus of
the field, a complex torus that is not algebraic, on which every coherent
sheaf has vanishing Chern character strictly between degree zero and the top
degree (\Cref{prop:weiltorussheaves}, \Cref{thm:p2false},
\Cref{thm:cmpurefalse}). With powers of the polarisation added the criterion
becomes one number: for a Chern character $N\omega+\sum_{k}c_{k}\theta^{k}$
and $n\ge3$ every object has $\dim\Ext^{2}(E,E)\ge(4+\rho)n^{2}-2n$, where
$\rho\le3$ is the rank of the Hankel matrix of the numbers $k!\,c_{k}$, and
an object attaining this bound meets the criterion
(\Cref{thm:p2primenumber}).%
\footnote{That is, $\rho$ is the rank of the $(2n-1)\times3$ matrix with
entries $\mu_{i+j}$, $0\le j\le2n-2$, $0\le i\le2$, where $\mu_{k}=k!\,c_{k}$
(\Cref{setup:p2prime}). A vector in its kernel is a linear recurrence of order
at most two among the $\mu_{k}$, and $\rho$ does not change when the character
is multiplied by $e^{t\theta}$.}
A flatness condition on the deformations of the object suffices in place of
injectivity (\Cref{prop:p2flat}). By descent and scalar extension the
statement is needed only for the split families of the CM fields of degree at
least four (\cite[Proposition~2.13]{paperIII}), and for one family it is equivalent to the
Lefschetz standard conjecture for the total space of the family over a curve
through a base point (\cite[Corollary~2.31]{paperIII}).

For a quartic CM field and $n=2$ we compute the rank exactly for every
shape. It takes fifteen values; the least, $80$, is taken only by the pure
shape, which the Weil tori exclude, so an object meeting the criterion has
$\dim\Ext^{2}(E,E)\ge88$ and a polynomial part that is exponential or linear
at each real place (\Cref{thm:quarticrank}, \Cref{cor:quarticleast}). The
Hodge locus of every such character is determined: it is the Weil tori of the
field only for the characters of $\theta^{4}$ shape, which are therefore
excluded at rank $88$, and the polarised family itself for every other
character of that rank (\Cref{prop:quarticlocus}, \Cref{cor:quarticlocus}).

\subsection{Secant objects}

The objects that a secant construction supplies are examined in
\Cref{sec:secantobjects}. On the split member $X\times\widehat{X}$ of
\cite[Section~14]{paperI} the secant plane is written down in every dimension,
and a secant sheaf has rank $a$, $c_{1}=bt$, $c_{2}=\tfrac{ad+b^{2}}{2}t^{2}$
and Bogomolov discriminant $\Nm_{K/\QQ}(b+a\sqrt{-d})\,t^{2}$
(\cite[Theorem~14.5]{paperI}, \cite[Theorem~14.13]{paperI}). No object whose Chern character
is an effective combination of exponentials of divisor classes contributes to
the Weil line, by a positivity argument with the norm form
(\Cref{thm:positivity}), and no object with a line bundle summand off the
polarisation line meets the criterion, so no sum of line bundles does in any
dimension (\cite[Theorem~9.20]{paperI}, \Cref{cor:noline}). Convolutions of line
bundles, their iterated cones, are cut down by the maps between the pieces.
Fewer than $2n$ pieces with nondegenerate differences never meet the
criterion (\Cref{thm:fewpieces}). On $E_{0}^{2n}$, $E_{0}=\CC/\ZZ[i]$, the
signed Chern characters of $2\cdot4^{n-1}$ line bundles add up to a real
Weil class (\Cref{prop:weilpieces}), and with three multiples of the
polarisation no convolution of these pieces meets the criterion at $n=3$
(\Cref{thm:noconvolution}), nor at $n=4$ when the shifts of the pieces lie
in six consecutive values (\Cref{cor:efoursix}). The weakened criterion
of \cite[Question~11.4]{Mar25b} reduces, for the transform of
$\cF\boxtimes\cF$, to the single identity
$\mathrm{At}(\cF)^{2}=-d\Theta^{2}\cdot\id$ for the sheaf $\cF$ on the
fourfold (\Cref{thm:factor}); a local computation at an isolated double point
of the support decides it in the negative for Markman's candidate in
dimension eight (\Cref{thm:candidatefails}), and at a quartic CM field a
curve on which a complex is locally free is tangent to at most one of the two
eigenplanes of the real multiplication, which decides Markman's explicit pair
\cite[Example~11.2.7]{Mar25c} in the same way (\Cref{thm:quarticlocal},
\Cref{cor:markmanquarticfails}).

The supports are constrained as well. A Cohen--Macaulay support carries no
local obstruction, while two smooth branches of codimension two meeting
transversally at a point carry one (\Cref{cor:cmdichotomy}). For a smooth support, a surface $S$ in the
fourfold, Grothendieck--Riemann--Roch fixes every numerical invariant, and the
Bogomolov--Miyaoka--Yau inequality gives $d\le b^{2}$ for the twist
$b\Theta$, which leaves only $d\in\{1,3,5,7\}$ for the twist $3\Theta$ of
Markman's candidate (\Cref{thm:smoothinvariants},
\Cref{cor:fourdiscriminants}). Complete
intersections are excluded for every $b$ and $d$ (\Cref{thm:noci}), two-term
complexes of powers of the polarisation at every rank, twist and
discriminant (\Cref{thm:nosplit}), and a Hilbert--Burch resolution of a
secant support has rank at least two (\Cref{prop:lowrankburch}); the secant
class is an integer combination of classes of line bundles exactly when
$d\equiv15$ or $23\pmod{24}$ (\Cref{thm:lattice}).


The paper is organised as follows. \Cref{sec:rankpoint} proves that one
semiregular object at a base point closes $\Sigma_{F}$ and draws the
consequences of injectivity for the object. \Cref{sec:polarised} shows that
the pure form of the criterion is never met and computes the polarised form
as a number. \Cref{sec:secantobjects} examines the secant objects of the
split member, the sums and convolutions of line bundles, the weakened
criterion for Markman's candidates and the supports. \Cref{sec:cmcriterion}
extends the exclusion of the pure form to every CM field and computes the
polarised criterion for quartic fields and the counts for sextic fields.

Every proof is written to be checked by hand. Where an argument rests on a
finite calculation, the calculation is given where it is used, as an explicit
identity, a determinant or a case analysis. The archive \cite{BMB26} holds
programs that repeat much of this arithmetic and a Lean~4 certificate of
part of it; no statement of the paper rests on them.
